\documentclass[../main.tex]{subfiles}
\begin{document}

\chapter*{Exam Guide --- Worked Example}
\addcontentsline{toc}{chapter}{Exam Guide --- Worked Example}
\markboth{Exam Guide --- Worked Example}{Exam Guide --- Worked Example}
\label{ch:exam}

% ---------------------------------------------------------------
\section*{Q1 — Referencing (\autoref{sec:referencing})}

\textbf{Author-date (APA)}
\begin{itemize}
  \item Author format: \texttt{Last, Initial.} (e.g.\ \texttt{Uriol Balbin, I., Bisagni, C., \ldots})
  \item Year in parentheses after authors: \texttt{(2020).}
  \item List ordered \textbf{alphabetically}
  \item In-text: \texttt{(Author, Year)} or \texttt{Author (Year)}
\end{itemize}

\textbf{Numbered (AIAA)}
\begin{itemize}
  \item Bracketed number before entry: \texttt{[1]}
  \item Author format: \texttt{Last, Initial.}
  \item Entry: \texttt{[1] Last, I. "Title," Journal, Vol. X, No. Y, Year, pp. X–Y.}
  \item List ordered by \textbf{order of appearance} in text
  \item In-text: \texttt{[1]}
\end{itemize}

Both styles require \texttt{Last, Initial} format. \texttt{Initial Last} and \texttt{First Last} are errors. Mixing formats within a list is an error.

% ---------------------------------------------------------------
\section*{Q2 — Research Question (\autoref{sec:objectives-rqs})}

\textbf{Criteria — RRAP}
\begin{itemize}
  \item \textbf{Relevant} — contributes to the objective
  \item \textbf{Researchable} — answerable through methods
  \item \textbf{Anchored} — tied to a specific system/context
  \item \textbf{Precise} — no vague language, no hedges (\emph{might, aspects of, could})
\end{itemize}

\textbf{Spotting the error (Q2a type)}
\begin{itemize}
  \item Vague words $\to$ Imprecise
  \item No system/context $\to$ Not anchored
  \item Yes/no formulation $\to$ Too closed
  \item Grammatical error and ``too short'' are almost never the answer
\end{itemize}

\textbf{Template:} ``What [specific attribute/relationship] of [specific concept] [affects/determines] [specific outcome] in [specific system/context]?''

\textbf{Examples}
\begin{itemize}
  \item[$\times$] ``What aspects related to thermal control might be considered important?''
  \item[$\checkmark$] ``What thermal control parameters most significantly affect structural integrity during spacecraft reentry?''
  \item[$\checkmark$] ``How does inlet pressure influence combustion efficiency in a scramjet operating at Mach 6?''
  \item[$\checkmark$] ``Which aerodynamic surface deflection angles minimize drag during the transition phase of a tilt-rotor tail-sitter MAV?''
\end{itemize}

Must be open-ended (\emph{What / How / Which / To what extent}), not yes/no.

% ---------------------------------------------------------------
\section*{Q3a — Research Objective (\autoref{sec:objectives-rqs})}

\textbf{Prescribed structure:} ``The research objective is \textbf{[contribution]} by \textbf{[method/approach]}.''

\textbf{Criteria — IUFC}
\begin{itemize}
  \item \textbf{Informative} — generic engineer can understand the project
  \item \textbf{Useful} — states what benefit the research provides
  \item \textbf{Feasible} — achievable within scope and resources
  \item \textbf{Clear} — precise about what the contribution is
\end{itemize}

Common error: writing a question (``How can X be improved\ldots'') instead of a declarative statement.

% ---------------------------------------------------------------
\section*{Q3b — Variables (\autoref{sec:language-research}, \autoref{sec:operationalization})}

A variable is any quantity that can change and is measured/controlled in the research. \textbf{Only select variables explicitly mentioned in the scenario text.} Do not import variables from background knowledge.

% ---------------------------------------------------------------
\section*{Q4 — Data Management (\autoref{sec:dmp})}

\textbf{3-2-1 rule:} 3 copies, on 2 different storage media, with 1 copy in a different physical location. If a data type has only one storage location $\to$ backup is missing; flag it.

\textbf{DMP principles:} Use open/non-proprietary formats. Keep an uncorrected original. Document your workflow. Use descriptive file names. Design the backup strategy from the start, not at the end.

% ---------------------------------------------------------------
\section*{Q5 — Research Design (\autoref{sec:language-research}, \autoref{sec:strategy-materials})}

\begin{description}
  \item[Hypothesis] Affirmative prediction that the proposed solution works. Mirrors the RO/RQ positively. Not the problem statement, not a comparison with prior methods.
  \item[Expertise] What you need to know. Only select domains directly required by the tasks listed.
  \item[Methods] What you do (analytical frameworks, simulation, testing, NDI/INDI controller). Knowledge areas are not methods.
  \item[Materials] Physical/software resources you procure (aircraft, computer, software, test facility). Methods and expertise are not materials.
\end{description}

% ---------------------------------------------------------------
\section*{Q6 — Sampling (\autoref{sec:sample-types})}

\begin{center}
\begin{tabular}{@{}lll@{}}
  \toprule
  \textbf{Type} & \textbf{Definition} & \textbf{Main limitation} \\
  \midrule
  Simple random & Equal probability for every element & Underrepresents rare sub-groups \\
  Convenience   & What is available                   & May not represent long-term or broader trends \\
  Purposive     & Target the information-rich sub-group & No comparison with the non-targeted group (no baseline) \\
  \bottomrule
\end{tabular}
\end{center}

Neither convenience nor purposive sampling can be assumed to represent the full population — must be stated explicitly.

% ---------------------------------------------------------------
\section*{Q7 — 4Cs Writing (\autoref{sec:writing-style})}

\begin{description}
  \item[Clear] Understood by intended audience; terms defined.
  \item[Comprehensive] All necessary information present.
  \item[Concise] No repetition, no padding.
  \item[Correct] Factually accurate, grammatically sound.
\end{description}

Check distractors against the actual text before selecting. If an option says ``X is not defined,'' verify it is actually undefined in the passage.

\end{document}
